\chapter{The Bot and the Web API}
\label{chap:interfaces}

Besides the terminal, PESAD can be consulted through a Telegram bot and through
a web API, which serves the expert-system page of the author's site. Both are
written in Python, and they drive the same Prolog engine and knowledge base
through a common engine layer, described first. In addition, the web API offers an
evaluation of the knowledge base from test results, without questions
(section~\ref{sec:evaluate}).

\section{The Engine Layer}

The engine layer (the Python package \code{pesad\_engine}) is the only interface
between Python and Prolog. It keeps one SWI-Prolog process, driven over the
Machine Query Interface~\cite{swipl-mqi}, and runs each consultation as a
worker thread on the Prolog side, with a thread-local working memory
(\code{bot\_session.pl}). Each thread executes the ordinary shell of
chapter~\ref{chap:usage} and stops at each question. Each question is returned to
Python as a \emph{screen}, which holds the console text of the turn, the kind of
answer expected (a command, a yes/no answer, a certainty, an option, a number)
and the structured description of the question produced by \code{ui\_topic/1}
(its identifier, its text in the chosen language, its options, the explanation
given to \emph{why}). Three modes of use are offered: (i) turn by turn,
delivering one raw answer at a time to an open session, as the bot does; (ii)
statelessly, as the \code{/step} endpoint does, with a function that starts a
fresh session, replays a list of answers after checking each one against the
options of its screen, returns the screen reached and discards the session;
(iii) without questions, on facts and test results given at once, as the
\code{/evaluate} endpoint does. Since the shell, the question protocols and the
explanations are those of the terminal, the front-ends differ only in
presentation, and a change to the knowledge base reaches all of them at once.

\section{The Telegram Bot}

The bot conducts a consultation turn by turn in a Telegram chat, the chat
identifier being the session identifier. It shows each question with buttons, instead of waiting for typed answers. Buttons cover the language, the uncertainty method (each with a one-line
description, e.g., ``the weakest answer decides'' for the certainty factors and
``small doubts on several answers add up to zero'' for the \L{}ukasiewicz
t-norm) and the options of the multiple-choice questions. A yes/no item becomes
a single question ``How true is it?'' with five steps, each of which stands for
the pair of answers that the terminal requests separately, i.e., the yes/no
answer and its certainty:
\begin{center}
\small
\begin{tabular}{@{}lcccccc@{}}
\toprule
Step & not at all & a little & partly & mostly & completely \\
\midrule
Answer and certainty & no, $1.0$ & yes, $0.25$ & yes, $0.5$ & yes, $0.75$ & yes, $1.0$ \\
Certainty of the item & $0$ & $0.25$ & $0.5$ & $0.75$ & $1$ \\
\bottomrule
\end{tabular}
\end{center}
A number in $[0,1]$ typed after each answer is tedious and, for most users,
hard to interpret, whereas a verbal scale of five steps is answered at once.
The price is that the certainties are restricted to five values, which is
rather coarse but (arguably) about as fine as a clinician's confidence in an answer can
be stated. A multiple-choice answer gets certainty $1$. A ``back''
button undoes the last answer, by restarting the session and replaying the
answers given before it. At the end the bot shows the ranked list of the investigation,
with the proof tree of each goal on request.

The bot is distributed as a Docker image, which the continuous integration of
the repository builds and publishes to the GitHub container registry at each
push to the main branch, after the validation suite has passed. On the virtual
machine that hosts it, a Docker Compose stack runs the bot together with
Watchtower, which polls the registry every two minutes and restarts the bot on
a new image when one is published. Hence, a change to the knowledge base
reaches the bot without manual steps.

\section{The Web API}

The web API is a FastAPI application, deployed on Google Cloud Run (the image is
built with Cloud Build). It is stateless: a request carries the whole
consultation, so that no session is kept on the server and any instance can
serve any request. Nothing is stored, and the answers are not logged.

\paragraph{Consultation.}
\code{POST /step} receives the language, the uncertainty method (\code{cf},
\code{fuzzy}, \code{fuzzy\_membership}, \code{lukasiewicz}, \code{hamacher} or
\code{einstein}) and the list of the answers given so far, as the values of the
options of the questions. The server starts a fresh session, replays the
answers and returns either the next question (with its identifier, text,
explanation and options) or the result, i.e., the goals with a certainty
greater than $0$, each with its proof tree. An answer that is not an option of
its question is rejected with status 400 and the index of the answer. Replaying
the whole consultation at every step makes the cost of a step grow with the
length of the consultation, which is the price of a server without state; the
client (a page of the site) needs nothing but the list of its answers. Requests are limited to
1000 answers and 64 KB.

\paragraph{Evaluation from test results.}
\code{POST /evaluate} runs the knowledge base without questions, on facts given
at once, and is described in the next section. \code{GET /findings} lists what
it accepts: the 30 findings with their admissible ranges, the items, the choices
and the numbers.

\section{Evaluation from Test Results}
\label{sec:evaluate}

A consultation by questions assumes that the clinician answers the items one by
one. In practice, much of the evidence for the neurodevelopmental disorders
comes as the scores of tests and structured interviews (an IQ, an
adaptive-behavior composite, achievement scores, symptom counts of the DIVA-5
interview for adult ADHD~\cite{kooij}, screening scales such as the
ASRS~\cite{asrs}, the AQ~\cite{aq} or the RAADS-R~\cite{raads}, the percentile
of the Movement ABC-2~\cite{mabc}), and the data are often incomplete. The
evaluation (\code{evaluate.pl}) takes a request holding the age (in years, or
the age band), a list of facts and, optionally, the goals or the chapters to
report (the neurodevelopmental chapter by default). A fact may be an item
answered true or false with a certainty, the choice of an option of a question,
a number (e.g., the duration in months), or a \emph{finding}. A finding is a
test result, which \code{findings.pl} converts into evidence for one or more
items; the degree of the evidence, in $[0,1]$, is the membership of the score
to the cutoff of the criterion, scaled by the weight of the source. For
instance, a full-scale IQ supports the standardized-test part of criterion A of
intellectual developmental disorder with degree $1$ up to $70$ and $0$ from
$76$, linearly in between. Similarly, a reading standard score supports
criterion B of specific learning disorder with degree $1$ up to $78$ and $0$
from $85$ (cf. the thresholds of the manual recalled in
section~\ref{sec:domain}).

\paragraph{Evidence model.}
The evidence for an item comes at three levels, each of which overrides the
ones below it: (i) a direct answer of the clinician; (ii) \emph{confirmatory}
evidence, from a structured interview, a standardized test or a record; (iii)
\emph{provisional} evidence, from a screening scale, which suggests a criterion
but cannot establish it. Evidence of the same level for the same item combines
by its maximum.

\paragraph{Bounds.}
The goals are evaluated twice. In the \emph{lower} run, a missing datum takes
the value that is against the goal (a missing symptom is absent, a missing
exclusion holds). In the \emph{upper} run, it takes the value in favor of the
goal, and the provisional evidence counts as missing. Polarity is reversed
under each premise required false, so that a missing ``better explained by''
counts against the goal in the lower run and in its favor in the upper one. Since the t-norms, the rule chaining (\ref{eq:consequent}) and the
criterion ask (\ref{eq:criterion}) are monotone in the certainty of each
premise (with the polarity accounted for), the two runs give a lower bound
$\underline{c}$ and an upper bound $\overline{c}$ on the certainty that the goal
would have if the missing data were answered, either way, with certainty $1$.
A goal (and, likewise, each of its criteria) is then reported as
\begin{equation}\label{eq:status}
\text{status} \;=\;
\begin{cases}
\text{present} & \text{if } \underline{c} \ge 0.5,\\
\text{absent} & \text{if } \overline{c} < 0.5,\\
\text{undetermined} & \text{otherwise.}
\end{cases}
\end{equation}
For each goal the response gives the status and the two bounds, and the
criteria with their status (\code{met}, \code{not\_met}, \code{unknown}, or
\code{not\_assessed} when the upper run did not reach them) and bounds. It also
gives the \emph{missing} data (the items, choices and numbers that the two runs
met without an answer, each with its question), the proof tree as text and as
nested nodes, the ICD-10-CM and ICD-11 codes, and the \emph{differentials},
i.e., the best goal, with status and bounds, of each disorder that the
differential-diagnosis section of the manual lists as to be ruled out. As an example,
take a 30-year-old with a DIVA-5 count of 7 inattention symptoms in adulthood
and 6 in childhood, symptoms in 3 settings and occupational impairment. For
ADHD, predominantly inattentive presentation, the evaluation gives the status
\code{undetermined} with bounds $[0, 0.753]$. Criteria B, C and D are met (with
lower bounds $0.8$, $0.95$ and $0.97$), whereas the duration, the
hyperactivity-impulsivity items (which must be \emph{absent} for this
presentation) and the exclusion E are missing, and they are listed as the data
to collect. Among the differentials, reactive attachment and disinhibited social
engagement disorder are absent (their criteria refer to a child, and the
knowledge base does not apply them to an adult), and the others are
undetermined. An invalid request (an unknown finding, a score out of the range
of the instrument, a missing age) is rejected with status 400 and the index of
the offending fact.
